{\small
\setlength{\tabcolsep}{3.5pt}
\begin{longtable}{@{}llrrr@{}}
\caption{Information choice on the same evaluated bit grid}\label{tab:sensor49}\\
\toprule
$T;p$ & Method & \shortstack{Bound-minimizing\\bits} & \shortstack{Net-selected\\bits} & \shortstack{Net-regret\\bound range} \\
\midrule
\endfirsthead
\toprule
$T;p$ & Method & \shortstack{Bound-minimizing\\bits} & \shortstack{Net-selected\\bits} & \shortstack{Net-regret\\bound range} \\
\midrule
\endhead
\bottomrule
\endfoot
2;1 & W & 16,14,12 & 6,6,4 & [0.002563,0.002596] \\
2;1 & F & 16,14,12 & 7,4,4 & [0.002387,0.002426] \\
2;4 & W & 16,14,12 & 6,6,4 & [0.002691,0.00272] \\
2;4 & F & 16,14,12 & 4,4,4 & [0.002564,0.002564] \\
3;1 & W & 16,14,12 & 7,4,4 & [0.003092,0.003107] \\
3;1 & F & 16,14,12 & 6,4,4 & [0.002992,0.003036] \\
3;4 & W & 16,15,13 & 4,4,4 & [0.003322,0.003322] \\
3;4 & F & 16,14,12 & 4,4,4 & [0.003237,0.003237] \\
\end{longtable}
}
\noindent{\footnotesize Bit triples correspond, in order, to prices $1/16384,1/4096,1/1024$ per coordinate-bit. The net selector minimizes the upper actual-net-cost interval. Its regret bound is against the best bit in the finite catalogue, not an estimated unconstrained optimum. Every price-specific value and all nonnegative-price envelopes are deposited.\par}
